\appendix

\section{Selberg Zeta Program: Computational Details}\label{app:selberg_details}

\subsection{Ramanujan Classification Protocol}\label{app:ramanujan_protocol}

For each simple Lie algebra $\g$ of rank $r$ and dimension $d_\g$, the adjoint representation defines a graph $\Gra_{\mathrm{ad}}$ as follows:

\begin{enumerate}
    \item Vertices: the $d_\g$ generators in the Cartan-Weyl basis.
    \item Edges: directed edges $a \to b$ whenever $[T_a, T_b] \neq 0$.
    \item Adjacency matrix: $A_{ab} = 1$ if $[T_a, T_b] \neq 0$, $0$ otherwise.
    \item Degree matrix: $Q_{ab} = \deg_a \delta_{ab}$.
\end{enumerate}

The Ramanujan ratio is computed as:

\begin{equation}
\Ramanujan = \frac{\max_{|\la_i| < q+1} |\la_i|}{2\sqrt{q}}
\end{equation}

where $\la_i$ are the adjacency eigenvalues and $q+1$ is the maximum degree. $\Ramanujan \leq 1$ indicates the Ramanujan property.

\subsection{Ihara Zeta Function Computation}\label{app:ihara_computation}

The Ihara zeta function is computed via the Hashimoto formula:

\begin{equation}
Z(u)^{-1} = (1-u^2)^{r-1} \det(I - uA + u^2 Q)
\end{equation}

where $r = d_\g - \mathrm{rank}(\g) + 1$ is the cyclomatic number. Poles are found by numerical root-finding of $\det(I - uA + u^2 Q) = 0$.

\subsection{Pole Migration Velocity}\label{app:pole_migration}

The pole migration velocity is defined as:

\begin{equation}
v_{\mathrm{mig}} = \frac{\Delta(\mathrm{inner\_pole\_ratio})}{\Delta h}
\end{equation}

where $\mathrm{inner\_pole\_ratio} = \min_i |u_i| / r_c$ with $r_c = 1/\sqrt{q}$. The velocity is computed across consecutive algebras in each family.

\subsection{Composite Order Parameter}\label{app:order_parameter}

The composite order parameter $\Phi(h)$ is constructed from three spectral observables:

\begin{equation}
\Phi = w_1 \cdot \Phi_{\mathrm{Ram}} + w_2 \cdot \Phi_{\mathrm{mix}} + w_3 \cdot \Phi_{\mathrm{pole}}
\end{equation}

where $\Phi_{\mathrm{Ram}}$ is the Ramanujan ratio, $\Phi_{\mathrm{mix}}$ is the mixing time, and $\Phi_{\mathrm{pole}}$ is the pole density. The weights $w_i$ are normalized so that $\sum w_i = 1$.

The fitted relation is:

\begin{equation}
\Phi(h) \approx 0.289 \cdot h^{0.557}
\end{equation}

with the critical value $\Phi_c = 0.31$ determining the Ramanujan phase boundary.

\subsection{Product Graph Composability}\label{app:composability_details}

For direct sum $\g_1 \oplus \g_2$, the adjoint graph is the disconnected union $\Gra_1 \cup \Gra_2$. The Ramanujan property of the product is tested under six definitions:

\begin{enumerate}
    \item D1: Spectral ratio of the product adjacency matrix.
    \item D2: Spectral ratio of the Hashimoto non-backtracking operator.
    \item D3: Ihara zeta pole concentration.
    \item D4: Mixing time ratio.
    \item D5: Spectral gap ratio.
    \item D6: Degree heterogeneity ratio.
\end{enumerate}

Under all six definitions, Ramanujan $\times$ Ramanujan $\to$ Ramanujan unanimously.

\subsection{Non-Compact Form Invariance}\label{app:noncompact_invariance}

For each compact algebra $\g_{\mathbb{R}}$ and its $12$ non-compact real forms $\g_{\mathbb{R}}^{(i)}$, the Ramanujan classification is tested. All $12/12$ forms preserve the Ramanujan classification of the compact counterpart. The classification is a \emph{Dynkin-diagram invariant}.
